\documentclass[12pt]{article}
\usepackage[margin=1in]{geometry}
\title{Technical Glossary}
\author{Documentation Team}
\begin{document}
\maketitle

\section{Core Concepts}

\begin{description}
\item[Glyph] The smallest unit of text in a PDF. Represents a single character
with position, font, and size metadata.

\item[Content Block] A group of related text lines that form a logical unit such
as a paragraph, heading, or list item.

\item[Font Catalog] A registry of all fonts used in a document, including the
identified body font and heading fonts.

\item[Reading Order] The sequence in which content blocks should be read to
reconstruct the logical flow of the document.

\item[Intermediate Representation] The structured data model that connects all
pipeline layers, containing blocks, their types, and metadata.
\end{description}

\section{Pipeline Stages}

\begin{description}
\item[Extraction] Layer 1. Raw PDF data retrieval using PyMuPDF.
\item[Assembly] Layer 2. Character-to-block construction.
\item[Layout Analysis] Layer 3. Spatial structure detection.
\item[Semantic Analysis] Layer 4. Content type classification.
\item[Table Extraction] Layer 5. Structured table recovery.
\item[Emission] Layer 6. Markdown text generation.
\end{description}

\section{File Formats}

\begin{description}
\item[PDF] Portable Document Format. Page-oriented, position-based rendering.
\item[Markdown] Lightweight semantic markup. Line-oriented, structure-based.
\item[YAML] YAML Ain't Markup Language. Used for frontmatter metadata.
\item[ZIP] Archive format for bundling Markdown output with extracted images.
\end{description}
\end{document}